% =========================================================================
% EPÍLOGO DO LIVRO: SÍNTESE E HORIZONTES DE PESQUISA
% =========================================================================
\chapter*{Epílogo: Síntese e Horizontes de Pesquisa}\label{sec:conclusoes}
\addcontentsline{toc}{chapter}{Epílogo: Síntese e Horizontes de Pesquisa}

Os resultados e implementações consolidados neste volume integram a teoria matemática da Otimização Combinatória à engenharia de software em C++23. O texto estruturou-se em quatro eixos:

\begin{enumerate}[leftmargin=1.5em, itemsep=0.4em]
    \item \textbf{Fundamentação Teórica:} Demonstração matemática dos teoremas estruturais de fluxo (Fluxo Máximo e Corte Mínimo, Integralidade, Decomposição de Fluxo, Teoremas de Menger, König e as Quatro Dualidades em grafos bipartidos; e, em custo mínimo, otimalidade por ciclos negativos, custos reduzidos, folga complementar e Total Unimodularidade da matriz de incidência).
    
    \item \textbf{Arquitetura em C++23:} Desenvolvimento de classes base abstratas e algoritmos especializados divididos em duas famílias:
    \begin{itemize}
        \item \textbf{Família \lstinline{FlowNetwork} (Fluxo Máximo):} \lstinline{FordFulkerson}, \lstinline{EdmondsKarp}, \lstinline{Dinic}, \lstinline{PushRelabel} (FIFO) e \lstinline{PushRelabelImproved} (Gap Heuristic), documentados nos Apêndices~\ref{ap:flownetwork} a~\ref{ap:push_relabel_improved}.
        \item \textbf{Família \lstinline{CostNetwork} (Custo Mínimo):} \lstinline{CycleCanceling}, \lstinline{SuccessiveShortest} (SPFA), \lstinline{SuccessiveShortestDijkstra} (Potenciais) e \lstinline{NetworkSimplex}, documentados nos Apêndices~\ref{ap:costnetwork} a~\ref{ap:network_simplex}.
    \end{itemize}
    
    \item \textbf{Reduções e Aplicações:} Resolução de problemas clássicos (\textit{Maximum Bipartite Matching}, \textit{Minimum Vertex Cover}, \textit{Maximum Independent Set}, caminhos disjuntos e corte mínimo) e aplicação em \textbf{Segmentação Interativa de Imagens} via corte mínimo em grafos (\textit{Graph Cuts}).
    
    \item \textbf{Avaliação Experimental:} Protocolo de medição de tempo de CPU sob isolamento em instâncias padronizadas da literatura (coleções DIMACS), validando a corretude das soluções por certificação cruzada.
\end{enumerate}

Como trabalhos futuros, destacam-se a avaliação empírica sobre grafos esparsos de grande porte, a implementação de versões paralelas em OpenMP para algoritmos de pré-fluxo e a exploração de regras de pivoteamento adicionais no \textit{Network Simplex} (como \textit{Candidate List Pricing}).


